% Gerado por artigo/analise/para_latex.py a partir de artigo/resultados/numeros.json. Não editar à mão.
\begin{table}[htbp]
\centering
\fontsize{9}{11}\selectfont
\setlength{\tabcolsep}{4pt}
\caption{Mudança de tendência e variação percentual anual das taxas de jogo patológico e do controle negativo e dos registros do grupo de referência na atenção primária à saúde (APS) e nos Centros de Atenção Psicossocial (CAPS). Brasil, 2018--2026.}
\label{tab:tabela2}
\begin{tabular}{>{\raggedright\arraybackslash}p{7.2cm}R{1.4cm}R{2.6cm}R{1.4cm}R{2.6cm}}
\toprule
Indicador & APS & IC95\% (APS) & CAPS & IC95\% (CAPS) \\ \midrule
Jogo patológico: mês da mudança de tendência & jun/2022 & abr/2021--dez/2022 & mai/2024 & jun/2023--jan/2025 \\
Jogo patológico: VPA antes da mudança (\%) & 4 & \textminus{}26; 15 & 18 & 8; 25 \\
Jogo patológico: VPA depois da mudança (\%) & 100 & 91; 110 & 181 & 107; 367 \\
Controle negativo: VPA do último segmento (\%) & 4 & 1; 14 & 35 & \textminus{}1; 61 \\
Registros da referência: VPA do último segmento (\%) & 19 & 18; 20 & 6 & 3; 12 \\
\bottomrule
\end{tabular}
\par\smallskip\parbox{\linewidth}{\raggedright APS: atenção primária à saúde; CAPS: Centros de Atenção Psicossocial; IC95\%: intervalo de 95\% de confiança; VPA: variação percentual anual. Jogo patológico e controle negativo em taxa por 10 mil registros da referência; registros da referência em número absoluto (crescimento da rede). Regressão segmentada (\emph{joinpoint}) binomial negativa, ajustada pelo mês do ano e pela pandemia de COVID-19; IC95\% por \emph{bootstrap} de blocos de seis meses (999 réplicas). Último segmento do controle negativo: desde mai/2020 na APS e desde jul/2025 nos CAPS (IC95\% do mês: fev/2020--jul/2025). Referência: sem mudança de tendência na APS; nos CAPS, queda de jun/2019 a jun/2020 e crescimento desde então. Jogo patológico: F63.0; controle negativo: F63.1--F63.3. Referência: F32, F33, F41, F43, F10 e F20. APS: fev/2018 a set/2026; CAPS: jan/2018 a jun/2026. Fontes: Sistema de Informação em Saúde para a Atenção Básica; Sistema de Informações Ambulatoriais do SUS, Registro das Ações Ambulatoriais de Saúde.}
\end{table}
